<<<<<<<< HEAD:Sections/6. System Analysis/5-storage.tex
\subsubsection{Dịch vụ lưu trữ}
========
% File: fileservice.tex

\subsubsection{Công nghệ khác}
>>>>>>>> origin/main:Sections/6. System Analysis/2-Technology solutions/Other Tech.tex

\paragraph{Amazon S3}

\subparagraph{Giới thiệu}
Amazon S3 (Simple Storage Service) là một dịch vụ lưu trữ đối tượng (object storage) của AWS, được thiết kế để lưu trữ và truy xuất lượng lớn dữ liệu một cách an toàn và hiệu quả. Trong hệ thống TravelEZ, S3 được sử dụng chuyên biệt cho việc lưu trữ các file media \textbf{(hình ảnh, video)} mà người dùng tải lên khi viết review hoặc đăng bài blog. Việc tách biệt lưu trữ file ra khỏi máy chủ ứng dụng giúp giảm tải cho backend và dễ dàng mở rộng dung lượng lưu trữ trong tương lai.

\subparagraph{Ưu điểm}
\begin{itemize}
    \item \textbf{Độ bền và tính sẵn sàng cao}: S3 được thiết kế với độ bền dữ liệu lên tới 99.999999999\%, đảm bảo hình ảnh và video của người dùng gần như không bao giờ bị mất.
    \item \textbf{Khả năng mở rộng gần như vô hạn}: Bạn có thể lưu trữ bao nhiêu file tùy thích mà không cần lo lắng về việc hết dung lượng ổ đĩa trên máy chủ.
    \item \textbf{Chi phí hiệu quả}: Mô hình trả phí theo dung lượng sử dụng (pay-as-you-go) rất tiết kiệm, đặc biệt khi so sánh với chi phí nâng cấp ổ cứng cho máy chủ.
    \item \textbf{Tối ưu tốc độ truy cập}: Dễ dàng tích hợp với mạng lưới phân phối nội dung (CDN) như Amazon CloudFront để tăng tốc độ tải hình ảnh/video cho người dùng trên toàn cầu.
    \item \textbf{Giảm tải cho Backend}: Máy chủ Spring Boot sẽ không phải xử lý việc truyền tải các file tĩnh nặng, giúp tập trung tài nguyên cho việc xử lý logic nghiệp vụ.
\end{itemize}

\subparagraph{So sánh Amazon S3 và các lựa chọn khác}
\clearpage
\begin{landscape}
\begin{table}[h!]
\centering
\setlength{\tabcolsep}{4pt}
\footnotesize
\begin{tabularx}{\linewidth}{|l|X|X|X|}
\hline
\textbf{Tiêu chí} & \textbf{Amazon S3} & \textbf{Tự lưu trên server backend} & \textbf{Cloudinary} \\
\hline
\textbf{Khả năng mở rộng} & \textbf{Rất cao:} Gần như vô hạn. & \textbf{Thấp:} Phụ thuộc vào dung lượng ổ cứng của server. & \textbf{Rất cao:} Tương tự S3. \\
\hline
\textbf{Độ bền dữ liệu} & \textbf{Rất cao:} Tự động sao lưu, nhân bản trên nhiều nơi. & \textbf{Thấp:} Phải tự thiết lập cơ chế backup phức tạp. Nguy cơ mất dữ liệu cao nếu server gặp sự cố. & \textbf{Rất cao:} Tương tự S3. \\
\hline
\textbf{Hiệu suất} & \textbf{Cao:} Tối ưu cho việc lưu và đọc file. Dễ tích hợp CDN. & \textbf{Trung bình:} Bị ảnh hưởng bởi tải của server. Có thể gây nghẽn khi nhiều người truy cập file. & \textbf{Rất cao:} Tích hợp sẵn CDN và các công cụ tối ưu hình ảnh/video tự động. \\
\hline
\textbf{Chi phí} & \textbf{Linh hoạt:} Trả theo dung lượng sử dụng, chi phí thấp. & \textbf{Chi phí ẩn:} Tốn tài nguyên CPU, RAM và băng thông của server chính. Tốn chi phí nâng cấp ổ cứng. & \textbf{Cao hơn:} Thường đắt hơn S3 nhưng đi kèm nhiều tính năng xử lý media. \\
\hline
\textbf{Độ phức tạp} & \textbf{Thấp:} Chỉ cần gọi API để tải lên/tải xuống. & \textbf{Trung bình:} Phải tự quản lý cấu trúc thư mục, phân quyền file, và xử lý backup. & \textbf{Thấp:} Tương tự S3, API thân thiện. \\
\hline
\end{tabularx}
\caption{So sánh Amazon S3 và các lựa chọn khác}
\end{table}
\end{landscape}
